\section{Simulation Verification}
\label{sec:simulation}

The proposed controllers can be evaluated through a planar fencing scenario
with $N$ double-integrator agents and one moving target.  The agents are
initialized in a collision-free configuration, and the prescribed
performance function is chosen either as the exponential profile
\eqref{eq:rho-exp} or as the prescribed-time profile \eqref{eq:rho-pt}.  The
same artificial-potential function is used in all tests so that the safety
margin $d$ and activation radius $\mu$ are identical across controllers.

The main performance index is the enclosure error
\[
e_c(t)=\norm{\bar p(t)-p_0(t)}.
\]
For a successful PPC fencing response, $e_c(t)$ must remain below the
prescribed envelope $\rho(t)$ for the entire simulation interval.  Since
$\dist(p,p_0)\le e_c(t)$, this also provides an upper bound on the distance
from the target to the agents' convex hull.  The safety performance is
checked by plotting
\[
d_{\min}(t)=\min_{i\neq j}\norm{p_i(t)-p_j(t)}
\]
and verifying $d_{\min}(t)>d$ for all $t$.

To highlight the effect of prescribed performance, the PPC controllers
should be compared with a baseline fencing controller that uses the same APF
term but does not impose the envelope $\rho(t)$.  The comparison should
report the maximum enclosure error, the settling time relative to a selected
tolerance, the minimum pairwise distance, and the final velocity mismatch
\[
e_v(t)=\max_{i\in\NN}\norm{v_i(t)-v_0(t)}.
\]
These metrics directly correspond to the desired properties P4, P2, and P3.

For the APF-filtered individual PPC algorithm in \S\ref{sec:result2}, an
additional diagnostic is useful:
\[
e_{\eta,\max}(t)=\max_{i\in\NN}\norm{\eta_i(t)}.
\]
The local PPC constraints require $e_{\eta,\max}(t)<\max_i\rho_i(t)$, while
the centroid bound follows from the averaged local envelopes in
\eqref{eq:rho-split}.  This diagnostic separates the local tracking behavior
from the induced centroid-level fencing performance.

For the adaptive unknown-acceleration controller in \S\ref{sec:result3}, the
simulation should also report the adaptive gains $\hat\kappa_i(t)$.  The
test should include bounded but unknown target maneuvers and bounded agent
disturbances, so that the controller is evaluated without using
$u_0+d_0$ or $d_i$ in the feedback law.  Successful behavior is indicated by
bounded adaptive gains, preservation of the PPC envelope, and maintenance of
the minimum pairwise distance above $d$.

\subsection{Result 3 numerical example}

The MATLAB script \texttt{simulation/sim\_result3.m} implements the adaptive
APF-filtered PPC controller of \S\ref{sec:result3}.  The target total
acceleration and the agent disturbances are generated as bounded sinusoidal
signals, but they are not used by the controller.  In the numerical
implementation, the discontinuous sign function in \eqref{eq:r3-ui} is
replaced by a saturation function with a small leakage term in the adaptive
gain update, as discussed in Result~3.

Figure~\ref{fig:result3-centroid} shows that the centroid-to-target distance
remains below the prescribed envelope throughout the simulation.  The final
centroid error is $1.4\times10^{-3}$, while the final envelope value is
$5.65\times10^{-1}$.

\begin{figure}[t]
    \centering
    \includegraphics[width=\columnwidth]{simulation/out/result3_centroid_ppc.pdf}
    \caption{Result 3: centroid-to-target distance and prescribed
    performance envelope.}
    \label{fig:result3-centroid}
\end{figure}

Figure~\ref{fig:result3-pairwise} plots all inter-agent distances.  The
minimum pairwise distance over the simulation is $1.0487$, which remains
above the safety threshold $d=0.65$.

\begin{figure}[t]
    \centering
    \includegraphics[width=\columnwidth]{simulation/out/result3_pairwise_distances.pdf}
    \caption{Result 3: pairwise inter-agent distances and the safety
    threshold.}
    \label{fig:result3-pairwise}
\end{figure}
